\chapter{Model-Checking Algorithms}
\label{ch:model-checking-algorithms}

A temporal specification becomes useful when a checker can either establish it or return a diagnostic that an engineer can replay.
For a finite protocol, reachability is the first useful checker: enumerate states, stop at a bad state, and retain the route that led there.
Branching-time checking extends this idea by computing sets of states that satisfy subformulas.
Symbolic methods change the representation of those sets, while abstraction changes which states the checker distinguishes.
Each change has a separate correctness obligation.

Throughout this chapter, $K=(S,I,R,AP,L)$ is a finite Kripke structure with the totalized transition relation of Section~\ref{sec:ch45-deadlocks}.
There are no fairness restrictions unless explicitly stated.
Write $V=|S|$ and $E=|R|$.
State identifiers are dense integers, adjacency lists enumerate all successors, and label tests take constant time.
These representation assumptions make the graph bounds below meaningful; generating successors from a program may have substantial additional cost.

\section{Reachability with a shortest witness}
\label{sec:ch46-reachability}

Figure~\ref{fig:ch46-work-graph} models a compilation job.
State $s_0$ dispatches it, $s_1$ waits for a resource and may retry, $s_2$ has prepared installation metadata, and $s_3$ has installed the code.
An alternative dispatch reaches $s_4$, an invalid allocation outcome, and then the explicit error state $s_5$.
We use $I=\{s_0\}$, $\mathit{installed}$ at $s_3$, and $\mathit{error}$ at $s_5$.
Normal completion and error have self-loops.
This is a deliberately faulty protocol: the checker should find the allocation-failure route even though a successful route also exists.

\begin{figure}[tbp]
\centering
\begin{tikzpicture}[node distance=15mm and 18mm]
\node[basic block] (s0) {$s_0$: dispatch};
\node[basic block,right=of s0] (s1) {$s_1$: wait};
\node[basic block,right=of s1] (s2) {$s_2$: prepared};
\node[basic block,right=of s2] (s3) {$s_3$: installed};
\node[basic block,below=of s1] (s4) {$s_4$: invalid};
\node[basic block,right=of s4] (s5) {$s_5$: error};
\draw[cfg edge] ([xshift=-8mm]s0.west) -- (s0);
\draw[cfg edge] (s0) -- (s1);
\draw[cfg edge] (s1) -- (s2);
\draw[cfg edge] (s2) -- (s3);
\draw[cfg edge] (s0) -- (s4);
\draw[cfg edge] (s4) -- (s5);
\draw[cfg edge] (s1) to[loop above] (s1);
\draw[cfg edge] (s3) to[loop above] (s3);
\draw[cfg edge] (s5) to[loop right] (s5);
\end{tikzpicture}

\caption{A successful installation route and an allocation-error route coexist in one finite model.}
\label{fig:ch46-work-graph}
\end{figure}

Let $B\subseteq S$ be the bad states, here $\{s_5\}$.
The invariant that no initial execution visits $B$ is equivalent to $\mathrm{Reach}(I)\cap B=\emptyset$.
Algorithm~\ref{alg:ch46-bfs} uses \term{breadth-first search} (BFS) and records one parent when a state is first discovered.
All initial states enter the queue before exploration, so the witness is shortest from the entire initial set, not merely from the first initial state considered.
The parent marker $\bot$ marks a root and is not a state identifier.

\begin{algorithm}[tbp]
\caption{Search for a bad state and retain a shortest transition witness.}
\label{alg:ch46-bfs}
\begin{pseudocode}
\Require finite graph with complete adjacency lists, initial set $I$, bad set $B$
\Ensure a shortest path from $I$ to $B$, or proof that no such path exists
\State $D\gets\emptyset$; $W\gets$ empty FIFO queue
\ForAll{$s\in I$}
  \State $D\gets D\cup\{s\}$; $\mathit{parent}[s]\gets\bot$; enqueue $s$ in $W$
\EndFor
\While{$W\neq\emptyset$}
  \State $s\gets$ dequeue from $W$
  \If{$s\in B$}
    \State \Return reverse of the parent chain from $s$ to $\bot$
  \EndIf
  \ForAll{$t\in\Succ(s)$}
    \If{$t\notin D$}
      \State $D\gets D\cup\{t\}$; $\mathit{parent}[t]\gets s$; enqueue $t$ in $W$
    \EndIf
  \EndFor
\EndWhile
\State \Return unreachable
\end{pseudocode}
\end{algorithm}

The discovery invariant is that every state in $D$ has a parent chain beginning in $I$ and following actual edges.
The FIFO queue processes nondecreasing distances, so first discovery gives a minimum edge count.
Marking a state when enqueuing it, rather than when removing it, prevents many predecessors from inserting duplicate queue entries.
Cycles and terminal self-loops therefore cause no special difficulty.
If the queue empties, $D$ contains $I$ and is closed under successors, which proves that no reachable bad state was missed.

For dense identifiers and adjacency lists, each state is enqueued once and each outgoing edge of an explored state is scanned once.
Worst-case time is $O(V+E)$, with $O(V)$ auxiliary memory for the queue, discovery flags, and parents, in addition to the $O(V+E)$ graph.
Reconstructing a returned path takes time proportional to its length, at most $V-1$ edges for a shortest nontrivial witness.
For an implicit graph, charge separately for successor generation, canonicalization, and state comparison.
Hash tables give expected constant-time lookup under suitable assumptions, not a representation-independent worst-case linear bound.

Table~\ref{tab:ch46-bfs-levels} computes the BFS layers with successors listed in numeric order.
The checker returns $s_0,s_4,s_5$, a two-transition witness.
The installation state at depth three is irrelevant to this violation.
The displayed layers describe a complete BFS traversal; an early-stopping implementation need not finish discovering the last layer before reporting the error.

\begin{table}[tbp]
\centering
\caption{BFS distances distinguish the shortest error route from the longer successful route.}
\label{tab:ch46-bfs-levels}
\begin{tabular}{cll}
\toprule
Distance & Newly reached states & Recorded parents \\
\midrule
0 & $\{s_0\}$ & $s_0:\bot$ \\
1 & $\{s_1,s_4\}$ & $s_1:s_0,\ s_4:s_0$ \\
2 & $\{s_2,s_5\}$ & $s_2:s_1,\ s_5:s_4$ \\
3 & $\{s_3\}$ & $s_3:s_2$ \\
\bottomrule
\end{tabular}
\end{table}

A witness should include enough decoded state to explain the failed condition.
For a cache model, pin ownership and reclamation status matter more than a raw state number.
If several actions connect the same pair of states, retain the generating action alongside the parent.
Otherwise the displayed state path may be valid while its claimed explanation names the wrong operation.
Chapter~\ref{ch:model-checker-case-study} provides a companion implementation of finite reachability and CTL checking.

\section{CTL as set computation}
\label{sec:ch46-ctl-sets}

For a CTL state formula $\varphi$, write
\[
 \llbracket\varphi\rrbracket_K=\{s\in S\mid K,s\models\varphi\}.
\]
We omit the subscript when the model is fixed.
Boolean cases are set operations: negation is complement relative to $S$, conjunction is intersection, and an atomic proposition is its label set.
Define the existential predecessor operation
\[
 \operatorname{Pre}_{\exists}(X)=\{s\in S\mid\exists t\in X:(s,t)\in R\}.
\]
Then $\llbracket\mathrm{EX}\varphi\rrbracket=\operatorname{Pre}_{\exists}(\llbracket\varphi\rrbracket)$ and $\llbracket\mathrm{AX}\varphi\rrbracket=S\setminus\operatorname{Pre}_{\exists}(S\setminus\llbracket\varphi\rrbracket)$.
Evaluating subformulas from the leaves upward turns temporal nesting into ordinary dependencies between set computations.

\subsection{Until is a least fixed point}
\label{sec:ch46-eu}

Let $P=\llbracket\varphi\rrbracket$ and $Q=\llbracket\psi\rrbracket$.
A witness for $\mathrm{E}[\varphi\,\mathrm{U}\,\psi]$ either starts in $Q$, or starts in $P$ and takes a step to a shorter witness.
Consequently,
\[
 \llbracket\mathrm{E}[\varphi\,\mathrm{U}\,\psi]\rrbracket
 =\mu X.\bigl(Q\cup(P\cap\operatorname{Pre}_{\exists}(X))\bigr).
\]
Here $\mu$ denotes the least fixed point under set inclusion.
Iterating from $\emptyset$ first yields $Q$; equivalently, start at $X_0=Q$ and repeatedly apply the displayed function.
Then $X_j$ contains exactly the states with an until witness reaching $Q$ in at most $j$ transitions.
Induction on that length proves both directions of the characterization.
Monotonicity and finiteness imply stabilization after at most $V$ strict state additions, although a round can add many states.

The least fixed point matters.
If $u\in P\setminus Q$ has only $u\to u$, it must not satisfy strong until.
Starting from all states can leave $u$ in a larger fixed point even though no target is ever reached.
Reaching $Q$ immediately, however, imposes no requirement that the current state also belong to $P$.
This zero-step case is a frequent implementation bug.

Algorithm~\ref{alg:ch46-eu-worklist} uses reverse edges instead of rescanning the graph in every round.
For each newly admitted predecessor it records a successor already known to reach $Q$.
These pointers lead to a target because they decrease discovery rank, even if the original graph contains cycles.

\begin{algorithm}[tbp]
\caption{Compute existential until by backward propagation from its target set.}
\label{alg:ch46-eu-worklist}
\begin{pseudocode}
\Require finite graph with predecessor lists; state sets $P$ and $Q$
\Ensure $X=\llbracket\mathrm{E}[P\,\mathrm{U}\,Q]\rrbracket$ and witness pointers outside $Q$
\State $X\gets Q$; $W\gets$ FIFO queue containing every state of $Q$
\While{$W\neq\emptyset$}
  \State $t\gets$ dequeue from $W$
  \ForAll{$s\in\Pred(t)$}
    \If{$s\in P$ and $s\notin X$}
      \State $X\gets X\cup\{s\}$; $\mathit{next}[s]\gets t$; enqueue $s$ in $W$
    \EndIf
  \EndFor
\EndWhile
\State \Return $X,\mathit{next}$
\end{pseudocode}
\end{algorithm}

Each admitted state is processed once, and each incoming edge of an admitted state is examined once.
The worst-case time is $O(V+E)$; auxiliary storage is $O(V)$ beyond the forward and reverse graph.
Constructing reverse adjacency lists also costs $O(V+E)$ time and space.
A full-round implementation of the mathematical recurrence can instead require $O(V(V+E))$ time.
The fixed-point equation alone therefore does not establish a linear implementation bound.

On Figure~\ref{fig:ch46-work-graph}, take $P=\{s_0,s_1,s_2,s_4\}$ and $Q=\{s_3\}$.
Table~\ref{tab:ch46-fixpoint-rounds} shows the full-round sequence for explanatory purposes.
State $s_4$ never enters: its only successor is the error state, which has no route to installation.
At $s_1$, the self-loop does not prevent existential until from holding; the witness chooses the edge to $s_2$.
The universal claim that every run eventually installs is nevertheless false.

\subsection{Always on some path is a greatest fixed point}
\label{sec:ch46-eg}

To satisfy $\mathrm{EG}\varphi$, a state must belong to $P=\llbracket\varphi\rrbracket$ and have a successor that can continue satisfying the same obligation.
Thus
\[
 \llbracket\mathrm{EG}\varphi\rrbracket
 =\nu X.\bigl(P\cap\operatorname{Pre}_{\exists}(X)\bigr),
\]
where $\nu$ denotes the greatest fixed point.
Start at $Y_0=P$ and repeatedly remove states having no successor in the current set.
The sequence descends and stabilizes.
Every retained state in the final set has a retained successor, so choosing successors indefinitely constructs a path entirely within $P$.
Conversely, no state on such a path can be removed: its next path state survives each round by induction.
This proves exactness.

For the same $P$ used in the until example, the right column of Table~\ref{tab:ch46-fixpoint-rounds} removes $s_2$ and $s_4$ immediately.
Their successors lie outside $P$.
States $s_0$ and $s_1$ survive because $s_0$ can move to $s_1$ and $s_1$ can repeat forever.
Starting this computation at $\emptyset$ would remain empty and miss both states.

\begin{table}[tbp]
\centering
\caption{Opposite initialization directions separate finite until witnesses from infinite stay-within-set witnesses.}
\label{tab:ch46-fixpoint-rounds}
\begin{tabular}{cll}
\toprule
Round & Until: $X_0=\{s_3\}$ & Always: $Y_0=P$ \\
\midrule
0 & $\{s_3\}$ & $\{s_0,s_1,s_2,s_4\}$ \\
1 & $\{s_2,s_3\}$ & $\{s_0,s_1\}$ \\
2 & $\{s_1,s_2,s_3\}$ & $\{s_0,s_1\}$ (fixed) \\
3 & $\{s_0,s_1,s_2,s_3\}$ & $\{s_0,s_1\}$ \\
4 & $\{s_0,s_1,s_2,s_3\}$ (fixed) & $\{s_0,s_1\}$ \\
\bottomrule
\end{tabular}
\end{table}

A linear-time implementation stores, for each state in $P$, the number of successors currently retained.
Initialize these counts by scanning edges, queue all zero-count states, and mark them removed when enqueuing.
When removing $t$, decrement the count of each still-retained predecessor for the edge to $t$; enqueue a predecessor when its count reaches zero.
Every state is removed at most once and every relevant edge is processed at most once, giving $O(V+E)$ time.
If adjacency lists contain duplicate edges, counts and decrements must use the same multiplicity convention.
After pruning, retain one successor within the result at each surviving state.
Following these pointers eventually repeats a state, yielding a finite prefix and a nonempty repeating cycle.

Equivalently, restrict the graph to $P$, find its strongly connected components, retain components with a directed cycle, and propagate backward within $P$.
A singleton component needs a self-loop to be cyclic.
The normal and error terminal loops inserted by totalization count as cycles; whether they satisfy $\mathrm{EG}\varphi$ depends on their labels.
For example, the error terminal satisfies $\mathrm{EG}\mathit{error}$ and fails $\mathrm{EG}\neg\mathit{error}$.

Boolean operations, $\mathrm{EX}$, existential until, and $\mathrm{EG}$ suffice for CTL.
Useful dualities include $\mathrm{AF}\varphi=\neg\mathrm{EG}\neg\varphi$ and $\mathrm{AG}\varphi=\neg\mathrm{EF}\neg\varphi$.
Universal until additionally excludes infinite postponement:
\[
 \mathrm{A}[\varphi\,\mathrm{U}\,\psi]
 =\neg\bigl(\mathrm{E}[\neg\psi\,\mathrm{U}\,(\neg\varphi\land\neg\psi)]
       \lor\mathrm{EG}\neg\psi\bigr).
\]
With linear graph routines and shared subformula results, checking a formula of size $|\varphi|$ takes $O(|\varphi|(V+E))$ time on the explicit model~\cite{baier-2008-model-checking}.
Storing every subformula's satisfaction set takes $O(|\varphi|V)$ bits in a bit-vector implementation, in addition to graph and worklist storage.
Some CTL failures require branching explanations; a single linear trace is not a universal diagnostic format for arbitrarily nested formulas.

\section{Symbolic sets and bounded paths}
\label{sec:ch46-symbolic}

\subsection{Boolean decision diagrams represent sets, not individual visits}
\label{sec:ch46-bdd}

Encode a state using Boolean variables $x=(x_1,\ldots,x_n)$.
A characteristic formula $X(x)$ denotes a state set, and $T(x,x')$ denotes the transition relation between current and next variables.
Let $D(x)$ describe valid state encodings; unused bit patterns must not become extra states.
The predecessor equation becomes
\[
 \operatorname{Pre}_{\exists}(X)(x)
 =D(x)\land\exists x'.\bigl(T(x,x')\land X(x')\bigr).
\]
The relation must constrain both ends to valid encodings, either internally or through explicit domain conjuncts.
Complement is $D\land\neg X$, not necessarily Boolean negation over every bit pattern.

A reduced ordered \term{binary decision diagram} (BDD) shares identical Boolean subfunctions under one fixed variable order.
Canonical representation for that order makes equality testing of normalized diagrams cheap and allows fixed-point iteration without enumerating each state.
For a two-bit state $(p,r)$ recording a pin and resident code, the safety set is $\neg p\lor r$.
It contains encodings $00,01,11$ and excludes $10$ when bits are ordered as $(p,r)$.
The set has a compact Boolean description despite containing several states.
The symbolic algorithm computes the same predecessor and fixed-point sets as the explicit algorithm; it changes their representation.

Compactness is not guaranteed.
Variable order can drastically change a BDD's size, and intermediate conjunctions or quantified results can be large even when the input diagrams are small.
Partitioning a transition relation and interleaving current and next variables are engineering choices to measure, not proofs of polynomial space.
A symbolic checker can avoid explicit enumeration on a favorable protocol and still exhaust memory on another with the same number of state bits~\cite{clarke-2018-handbook}.

\subsection{SAT unrolling asks a bounded reachability question}
\label{sec:ch46-bmc}

For a bad-state predicate $B$, an initial predicate $I$, and a transition formula $T$, define
\[
 \Phi_k=I(x_0)\land
       \bigwedge_{i=0}^{k-1}T(x_i,x_{i+1})\land
       \bigvee_{i=0}^{k}B(x_i).
\]
Each $x_i$ is a fresh copy of the state bits.
A satisfying assignment gives a length-$k$ execution with a bad state at some position at most $k$.
The formula must include the chosen arithmetic semantics, valid encodings, and all initial constraints.
In particular, mathematical integer arithmetic is not a substitute for wrapping machine arithmetic when overflow affects the protocol.

Because the transition relation is total, any shorter bad-reaching prefix extends to length $k$.
Thus this formula detects every bad state reachable within $k$ transitions.
On the six-state job model, $\Phi_0$ and $\Phi_1$ are unsatisfiable, while $\Phi_2$ is satisfiable with $s_0,s_4,s_5$.
Those first two UNSAT answers are correct and still do not establish safety.
They exclude only witnesses within their respective bounds.

For an explicitly known finite model with $V$ states, every reachable state has a simple shortest route of at most $V-1$ edges from some initial state.
Consequently, an exact UNSAT result at $k=V-1$ for this invariant query is complete.
This is a completeness argument tied to a particular state bound and property, not a general interpretation of bounded UNSAT.
An encoded model with $n$ bits may have up to $2^n$ states, making that bound unusable in practice.
Inductive invariants, diameter arguments, and other complete procedures can supply stronger certificates, but require their own checks.

\subsection{A lasso must satisfy the property automaton's acceptance condition}
\label{sec:ch46-lassos}

The retry counterexample $s_0s_1^\omega$ is a \term{lasso}: a finite prefix followed by a nonempty closed walk repeated forever.
For finite-state LTL checking, translate the negated property into an automaton over infinite label sequences and explore its product with the model.
A counterexample exists exactly when an accepting infinite run exists in that product; if one exists, an ultimately periodic witness can be chosen~\cite{baier-2008-model-checking}.
The repeated object is a product state containing both the system state and the automaton state.
Repetition of a system state alone does not establish that the pending temporal obligations repeat compatibly.

For ordinary B\"uchi acceptance, the repeated part must visit an accepting state.
For generalized B\"uchi acceptance, it must visit every required acceptance set.
Additional fairness requirements must also be satisfied.
A strongly connected region containing representatives of several required sets may supply a closed walk visiting them all, but no single simple cycle need do so.
Conversely, an arbitrary reachable cycle can be nonaccepting and prove nothing about the queried formula.

In a bounded lasso encoding with states $z_0,\ldots,z_k$, one convention chooses $\ell\in\{0,\ldots,k\}$ and requires a closing transition $z_k\to z_\ell$.
The repeated positions are $z_\ell,\ldots,z_k$; for B\"uchi acceptance at least one must be accepting.
Choosing $\ell=k$ requires a genuine self-loop.
An alternative convention repeats an equal endpoint and uses different indices; mixing the conventions can admit a zero-edge \enquote{cycle} that is not an infinite execution.
Finite bad prefixes suffice for invariants, while general temporal counterexamples need the appropriate acceptance test.
Neither all cycles nor all infinite paths are fair accepting witnesses.

\section{Abstraction and spurious traces}
\label{sec:ch46-abstraction}

Abstraction can make a protocol finite or make a finite protocol tractable.
It also introduces choices that may not compose into a concrete execution.
The soundness pattern resembles the over-approximate analyses of Chapter~\ref{ch:ch11-abstract-interpretation}, but now the abstraction must cover transitions as well as state facts~\cite{cousot-1977-abstract}.

Let $C$ and $A$ be concrete and abstract structures, and let $H\subseteq S_C\times S_A$ be a \term{forward simulation}.
Every concrete initial state must relate to an abstract initial state.
For every $(c,a)\in H$, labels must agree on the propositions used in the property, and every concrete transition $c\to c'$ must have a matching abstract transition $a\to a'$ with $(c',a')\in H$.
The abstract system may have additional transitions.
These conditions give a matching abstract path for every concrete path, with the same relevant labels.
Therefore universal LTL truth transfers from the abstract system to the concrete one, under the common infinite-path semantics.
They also preserve the universal fragment of CTL in the corresponding direction, but not arbitrary mixed existential and universal formulas~\cite{baier-2008-model-checking}.

This particular statement uses one abstract step per concrete step and exact agreement on relevant labels.
If states are merged despite differing labels, one needs a sound treatment of uncertain propositions rather than choosing convenient truth values.
If transitions are collapsed, next-time properties require special care.
If checking only fair abstract paths, every relevant fair concrete path must have a fair abstract match; ordinary simulation alone does not provide that extra condition.

Figure~\ref{fig:ch46-abstraction} gives a small counterexample to trusting an abstract error trace.
In the concrete pass protocol, $i\to u\to d$, and $d$ repeats forever.
State $w$ is unreachable from $i$, but has $w\to e$, and $e$ repeats forever.
States $u$ and $w$ share the visible phase \enquote{pending}; their hidden mode bits are 0 and 1, respectively.
Only $e$ has the proposition $\mathit{error}$.
Merge $u$ and $w$ into abstract state $A$, retaining other states as singleton blocks $I,D,E$.
An abstract edge exists whenever any pair of concrete representatives supplies it.
Then $I\to A\to E$ is an abstract error trace: its first edge uses $u$ and its second uses $w$.
There is no concrete execution that switches representatives this way.

\begin{figure}[tbp]
\centering
\begin{tikzpicture}[node distance=12mm and 15mm]
\node[basic block] (i) {$i$};
\node[basic block,right=of i] (u) {$u$: mode 0};
\node[basic block,right=of u] (d) {$d$: done};
\node[basic block,below=of u] (w) {$w$: mode 1};
\node[basic block,right=of w] (e) {$e$: error};
\draw[cfg edge] ([xshift=-8mm]i.west) -- (i);
\draw[cfg edge] (i) -- (u);
\draw[cfg edge] (u) -- (d);
\draw[cfg edge] (w) -- (e);
\draw[cfg edge] (d) to[loop right] (d);
\draw[cfg edge] (e) to[loop right] (e);
\node[draw,dashed,fit=(u)(w),inner sep=4pt,label={[font=\small]above:merge}] {};
\node[basic block,below=27mm of i] (ai) {$I$};
\node[basic block,right=of ai] (aa) {$A=\{u,w\}$};
\node[basic block,right=of aa] (ad) {$D$};
\node[basic block,below=of ad] (ae) {$E$};
\draw[cfg edge] ([xshift=-8mm]ai.west) -- (ai);
\draw[cfg edge] (ai) -- (aa);
\draw[cfg edge] (aa) -- (ad);
\draw[cfg edge] (aa) -- (ae);
\draw[cfg edge] (ad) to[loop right] (ad);
\draw[cfg edge] (ae) to[loop right] (ae);
\end{tikzpicture}

\caption{Merging pending states splices edges from incompatible concrete histories into an abstract error trace.}
\label{fig:ch46-abstraction}
\end{figure}

To validate an abstract trace $a_0,\ldots,a_k$, propagate concrete sets along the whole prefix:
\[
 C_0=I_C\cap\gamma(a_0),\qquad
 C_{j+1}=\operatorname{Post}_C(C_j)\cap\gamma(a_{j+1}).
\]
Here $\gamma(a)$ is the concrete block represented by $a$, and $\operatorname{Post}_C$ takes concrete successors.
For an abstract bad endpoint that is only a possible error, intersect the final set with the actual concrete bad set as well.
Table~\ref{tab:ch46-spurious} computes the sets for $I,A,E$.
Checking only that each adjacent pair has some concrete transition would incorrectly accept the trace.
The shared intermediate state must satisfy both steps in one consistent execution.

\begin{table}[tbp]
\centering
\caption{Prefix propagation rejects the abstract error trace when the concrete representative cannot continue.}
\label{tab:ch46-spurious}
\begin{tabular}{clll}
\toprule
Position & Abstract state & Concretization & Feasible prefix endpoint set \\
\midrule
0 & $I$ & $\{i\}$ & $\{i\}$ \\
1 & $A$ & $\{u,w\}$ & $\{u\}$ \\
2 & $E$ & $\{e\}$ & $\emptyset$ \\
\bottomrule
\end{tabular}
\end{table}

\term{Counterexample-guided abstraction refinement} (CEGAR) alternates abstract checking, concrete trace validation, and refinement of the distinctions responsible for a spurious result.
Here splitting $A$ by the hidden mode bit separates $u$ from $w$.
The refined reachable graph is $I\to U\to D$, with the terminal loop at $D$; $W\to E$ remains present but unreachable.
The abstract error disappears without deleting a real concrete transition.
Deleting the edge $A\to E$ without splitting or another sound argument would lose the behavior of $w$ and could invalidate the simulation contract.

For a finite concrete graph, repeated proper partition refinement can eventually reach singleton blocks, so this exact refinement scheme terminates if it always splits a block and never merges it again.
That statement is not a termination guarantee for general program CEGAR over unbounded stores.
Nor does feasibility of every finite unfolding of a proposed abstract loop by itself provide a concrete infinite fair witness under arbitrary abstraction.
Temporal counterexamples require validation of their looping and acceptance obligations as well as their finite prefix.

In a compiler workflow, a trace validator should distinguish a real protocol failure, an infeasible abstraction artifact, and an inconclusive solver result.
Only the first justifies reporting a concrete bug.
Only a completed sound check, with its modeling assumptions recorded, justifies the corresponding proof claim.
Reachability, temporal fixed points, symbolic encodings, and refinement are complementary components of that evidence chain.

\section{Exercises}
\label{sec:ch46-exercises}

\begin{enumerate}
\item\label{ex:ch46-multiple-initial}
On the six-state job model, change the initial set to $\{s_0,s_4\}$.
Find a shortest error witness and its length.
Explain why running BFS from $s_0$ to completion before considering $s_4$ does not implement multi-source shortest search.

\item\label{ex:ch46-until-zero}
Use the job model with $P=\{s_0,s_1\}$ and $Q=\{s_2,s_3\}$.
Compute $\mathrm{E}[P\,\mathrm{U}\,Q]$ by rounds starting at $Q$.
Which returned states do not belong to $P$, and why must they remain in the result?

\item\label{ex:ch46-eg-pruning}
Compute $\mathrm{EG}\neg\mathit{error}$ and $\mathrm{AF}\mathit{installed}$ at every state of the job model without fairness.
Give a nonterminating path explaining the answer at $s_1$.
Then remove only the edge $s_1\to s_1$ and recompute $\mathrm{AF}\mathit{installed}$.

\item\label{ex:ch46-bounded-unsat}
A total graph is a chain $t_0\to t_1\to t_2\to t_3\to t_4$ followed by $t_4\to t_4$.
Only $t_0$ is initial and only $t_4$ is bad.
For which nonnegative bounds $k$ is $\Phi_k$ satisfiable?
What would an UNSAT result at bound three justify?

\item\label{ex:ch46-accepting-walk}
A reachable strongly connected graph has edges $h\to u$, $u\to h$, $h\to v$, and $v\to h$.
Generalized B\"uchi acceptance requires infinitely many visits to both $\{u\}$ and $\{v\}$.
Find an accepting closed walk and explain why reporting either simple cycle alone would fail the acceptance test.

\item\label{ex:ch46-refinement}
For Figure~\ref{fig:ch46-abstraction}, validate the trace $I,A,D$ and then add $i\to w$ to the concrete graph.
Revalidate $I,A,E$ after that change.
Does splitting the pending block still eliminate the error, or does it expose a real witness?
\end{enumerate}

\section{Further reading}
\label{sec:ch46-reading}

Baier and Katoen give correctness arguments for explicit CTL algorithms, automata-based LTL checking, and simulation-based abstraction~\cite{baier-2008-model-checking}.
Huth and Ryan provide a set-based account of CTL fixed points that is useful when implementing the recurrences directly~\cite{huth-2004-logic}.
The \emph{Handbook of Model Checking} covers symbolic, SAT-based, and abstraction-refinement methods in greater implementation depth~\cite{clarke-2018-handbook}.
Cousot and Cousot provide the abstract-interpretation foundation for reasoning about sound approximations rather than treating merged states as exact executions~\cite{cousot-1977-abstract}.
